% LSPV-4 REPORT::DISTEQUIV1::2026-10-04T19:49:00-07:00
\documentclass[11pt]{article}
\usepackage[margin=.88in]{geometry}
\usepackage{amsmath,amssymb}
\begin{document}
\section*{LSPV-4 DISTEQUIV1: hostile distributional equivalence on an exact slice}
\textbf{GLOBAL STATUS: OPEN.} P versus NP remains open. No load-bearing implication is closed.

\section*{CORPUS REVIEWED}
Current repository governance, CODEX, the pnp-research skill, root main.tex, verified MAIN1854 and MAIN1625, TASK-001 through TASK-008 as surfaced by repository search, SLICEGATE1, SLICEPAD1, SLICEKW1, SLICEDEPTH2, SLICENEXT1, SLICELIFT2, DENSITYPOSFAIL1, SLICEDENSITY2, PAIRCAMOUFLAGE2, GROK H1927, RECURCLOSE1, SELECTCHARGE1, MULTICHANNEL1, DELAYOR1, RETURNSQUEEZE1, SOURCESTOP1, PATHSTRAT1, and OFFPATH1 were reviewed or reconciled through current verified artifacts. The newest verified MAIN found was MAIN1854. No summary-only failed report is promoted.

\section*{Target and quantifiers}
Fix a layer \(L_{N,t}=\{x\in\{0,1\}^N:|x|=t\}\) and arbitrary distributions \(\mu_0,\mu_1\) supported on that same layer.

\section*{KNOWN THEOREM used}
Valiant's slice theorem states that every Boolean function specified on one Hamming slice has monotone circuit complexity at most a constant factor above its unrestricted circuit complexity plus \(O(N(\log N)^2)\). Thus an unrestricted circuit \(C\) considered only on \(L_{N,t}\) can be replaced by a polynomially related monotone circuit \(M_C\) agreeing on every point of the layer.

\section*{PROVED/EXACT: DISTEQUIV1}
Let
\[
\operatorname{Adv}_{\mathcal C}(\mu_0,\mu_1)
=
\sup_{C\in\mathcal C}
\left|\Pr_{\mu_1}[C(x)=1]-\Pr_{\mu_0}[C(x)=1]\right|.
\]
For every polynomial size bound \(s(N)\), there is a polynomial \(p\) such that
\[
\boxed{
\operatorname{Adv}_{\mathrm{Circ}[s]}(\mu_0,\mu_1)
\le
\operatorname{Adv}_{\mathrm{MonCirc}[p(s,N)]}(\mu_0,\mu_1)
}
\]
and trivially
\[
\operatorname{Adv}_{\mathrm{MonCirc}[s]}(\mu_0,\mu_1)
\le
\operatorname{Adv}_{\mathrm{Circ}[s]}(\mu_0,\mu_1).
\]
Hence polynomial-size distributional distinguishability on one exact layer is the same for unrestricted and monotone circuits up to polynomial overhead.

\textbf{Proof.}
Take any unrestricted \(C\) of size at most \(s\). Restrict its Boolean function to \(L_{N,t}\). Valiant supplies a monotone circuit \(M_C\) of size \(p(s,N)\) agreeing with \(C\) at every point of the layer. Since both distributions have support in the layer, their acceptance probabilities under \(C\) and \(M_C\) are identical. Taking the supremum gives the first inequality. The reverse inclusion is immediate. \(\square\)

\section*{Hostile consequence for TASK-008}
Low-order camouflage is only a necessary filter. Matching all fixed-order degree, sign, overlap, codegree, and core statistics does not by itself approach the load-bearing lower bound. To prove a distributional lower bound against all polynomial-size monotone circuits on the exact slice, one must in effect prove indistinguishability against all polynomial-size unrestricted circuits on that same slice. Any polynomial-size unrestricted distinguisher discovered later becomes a polynomial-size monotone distinguisher with exactly the same layer advantage.

Therefore PAIRCAMOUFLAGE2 and H1927 are correctly finite hostile tests, not positive lower-bound evidence. H1927's observation that the apparent cubic signal largely disappears after pair-score conditioning is useful experimental hygiene, but even perfect fixed-order camouflage leaves DISTEQUIV1 untouched.

A positive distribution whose support membership lies in P/poly is especially fatal: its membership circuit gives a perfect unrestricted separator, hence a polynomial-size monotone separator on the layer. This recovers the P/poly positive-family guard in DENSITYPOSFAIL1.

\section*{Natural-Proofs control}
DISTEQUIV1 does not itself invoke Razborov--Rudich and does not prove that every successful TASK-008 argument is natural. It shows why monotone-first does not create a distributional loophole on a slice. If a proposed invariant is constructive and large in the Razborov--Rudich sense, the Natural-Proofs barrier must be audited separately. A target-specific non-large argument is not excluded.

\section*{Boundary-Squeeze audit}
RECURCLOSE1 and MULTICHANNEL1 remain exact recurrence accounting under genuine fractional-channel hypotheses. SELECTCHARGE1 and DELAYOR1 still expose current-scale channels. DISTEQUIV1 supplies no fractional shrinker. Williams' Circuit-SAT self-improvement is model-specific and does not furnish the missing source-matched OR-preserving selector for SU/vMS, CCSPC, or exact 3-SAT. DELAY, INTERP, RECOMP, semantic-oracle, and algebraic-oracle controls leave Boundary-Squeeze conditional.

\section*{Structural and causal acquisition audit}
Nothing in the slice theorem forces an arbitrary SAT algorithm to acquire SU/vMS product-span data or counterfactual CCSPC information. Compact representations remain an escape from representation-size arguments; semantic and algebraic oracle replacement continue to expose claims based only on extensional answers or algebraic collisions. SOURCESTOP1, PATHSTRAT1, and OFFPATH1 remain unchanged. DISTEQUIV1 supplies no arbitrary-algorithm compiler.

\section*{DUPLICATION / CONTRADICTION AUDIT}
PAIRCAMOUFLAGE2 states the point for any one polynomial-size pair-window classifier. DENSITYPOSFAIL1 states the perfect-separator special case for a P/poly-recognizable positive family. DISTEQUIV1 generalizes this to every unrestricted distinguisher and arbitrary same-layer distribution pair. It does not contradict SLICEGATE1, SLICEDEPTH2, or SLICELIFT2.

\section*{WHAT CHANGED RELATIVE TO ROOT MAIN}
The sharper statement is: on a fixed slice, distributional indistinguishability against polynomial-size monotone circuits is equivalent up to polynomial overhead to indistinguishability against polynomial-size unrestricted circuits. Thus fixed-order camouflage cannot be promoted as progress toward the general monotone lower bound without a theorem extending it to a growing circuit class.

\section*{ADOPTED / REJECTED / RETIRED}
\textbf{Adopted:} Valiant slice transfer, SLICEGATE1, the DAG/general monotone target, PAIRCAMOUFLAGE2 and H1927 only at finite computational scopes, RECURCLOSE1, SELECTCHARGE1, MULTICHANNEL1, and standing acquisition guards.

\textbf{Rejected as sufficient:} matching any fixed list of low-order local statistics; pair camouflage; absence of a detected order-three signal; large deletion-minimal cores found by heuristic search.

\textbf{Retired as a generic escape:} a distribution hard for polynomial-size monotone circuits on an exact slice while remaining distinguishable by polynomial-size unrestricted circuits on that slice.

\section*{Primary-literature transfer audit}
Valiant's 1986 slice theorem is a directly transferable theorem and is the sole theorem needed for DISTEQUIV1. Hrube\v{s}--Pudl\'ak's random-CNF monotone-real lower bounds remain for growing width \(k=\omega(1)\), not fixed 3-CNF. Williams' 2024 self-improvement theorem is a model-specific analogy for Boundary-Squeeze but supplies no source bridge. Razborov--Rudich is a barrier control only when constructivity, largeness, and usefulness are established.

\section*{Single next target}
For TASK-008, stop treating additional fixed-order camouflage as evidence of general hardness. Either prove a target-specific indistinguishability theorem against a genuinely expanding circuit class, with exact class and size bound, or find a structural monotone-DAG lower-bound invariant that is not merely a polynomial-size unrestricted slice predicate. Attack any proposed invariant with unrestricted polynomial-size classifiers before sunflower or DAG-KW promotion.

\section*{FINAL STATUS}
DISTEQUIV1 is PROVED/EXACT as a corollary of Valiant's KNOWN THEOREM. It is a hostile equivalence, not a lower bound. All load-bearing P versus NP arrows remain OPEN.

\section*{References}
L. G. Valiant, Negation is Powerless for Boolean Slice Functions, SIAM Journal on Computing 15(2):531--535, 1986.

A. A. Razborov and S. Rudich, Natural Proofs, Journal of Computer and System Sciences 55(1):24--35, 1997; ECCC TR94-010.

P. Hrube\v{s} and P. Pudl\'ak, Random Formulas, Monotone Circuits, and Interpolation, FOCS 2017.

R. Ryan Williams, Self-Improvement for Circuit-Analysis Problems, STOC 2024.

% END_OF_VERIFIED_REPORT::LSPV-4::DISTEQUIV1
\end{document}
